\documentclass[10pt,a4paper]{amsart}
\usepackage[margin=19mm]{geometry}
\usepackage{amsmath,amssymb,mathtools}
\usepackage[scr=rsfs,cal=euler]{mathalfa}
\usepackage{xfrac}
\usepackage[unicode,hidelinks,bookmarks=false]{hyperref}
\newcommand{\ie}{i.e.\ }
\newcommand{\cf}{cf.\ }
\newcommand{\resp}{respectively\ }
\newcommand{\Subsection}{\subsection}
\providecommand{\nobreakdash}{\nobreak\hbox{-}}
\setlength{\emergencystretch}{2em}
\multlinegap0pt
\usepackage{siunitx}
\usepackage{siunitx}
\usepackage{siunitx}
\usepackage{siunitx}
\usepackage[shortlabels]{enumitem}
\usepackage{amssymb}
\usepackage{mathtools}
\usepackage{bm}
\usepackage{siunitx}
\usepackage{tikz}
\usepackage{float}
\newtheorem{restatable}{Restatable}
\newtheorem{definition}{Definition}
\newtheorem{theorem}{Theorem}
\newcommand{\matr}[1]{\bm{#1}}        % Bold Matrices
\newcommand{\tr}{\operatorname{tr}}
\title{Data Collaboration Analysis with Orthonormal Basis Selection and Alignment}
\author{Keiyu Nosaka, Yamato Suetake, Yuichi Takano, Akiko Yoshise}
\date{Source: arXiv:2403.02780v9 (CC BY 4.0)}
\begin{document}
\maketitle

\noindent\emph{Source and licence.} Data Collaboration Analysis with Orthonormal Basis Selection and Alignment. Version arXiv:2403.02780v9 (CC BY 4.0), distributed by arXiv under the Creative Commons Attribution 4.0 International licence, http://creativecommons.org/licenses/by/4.0/, as recorded in the arXiv licence metadata for this exact version and shown on the abstract page https://arxiv.org/abs/2403.02780v9. Full licence notice: \texttt{licenses/arxiv-2403.02780-CC-BY-4.0.txt}.\par\smallskip

\noindent\textit{Editorial scope.} Counted unit: the article's theorem concordance (source0.tex lines 1332--1344). The article's complete original proof of it is delivered with it (source0.tex lines 1346--1385, one \texttt{proof} environment of 40 lines). No mathematics is written, repaired or re-derived: the statement and the proof are the article's own lines, in the article's own notation, and no retained line is substituted or re-flowed.  Retained uncounted as the source-local context the proof needs: lines 1227--1235; lines 1276--1280; lines 1323--1328.  Nothing was omitted from any retained span, so the package carries no omission record.  No retained line contains a citation, so the wrapper carries no bibliography and no bibliography entry.  No retained line contains a figure environment, an image-inclusion command or drawing code, so no image is delivered and the package declares no asset dependency.  Editorial formatting only, in addition: an AMS \texttt{amsart} preamble that restates the article's own declarations and macros used by the retained lines, namely the environments \texttt{restatable}, \texttt{definition}, \texttt{theorem} on separate counters, together with the standard packages those lines need.  The retained environments print in this document as Restatable 1, Definition 1, Theorem 1 in order of appearance, whereas the article prints the same items with the numbering of its own full document.  The complete native source of the article as distributed in the arXiv e-print for this exact version is bundled unchanged as \texttt{sources/155569/source0.tex}.
\begin{restatable}{assumption}{orthogonalassumptions}
    \label{assumptions: ODC}
    We impose the following conditions on the anchor dataset \( \matr{A} \in \mathbb{R}^{a \times m} \) and the secret bases \( \matr{F}_i \in \mathbb{R}^{m \times \ell} \; (m \geq \ell) \), for all \(i \in [c]\):
    \begin{enumerate}[label=\arabic*)]
        \item \( \mathrm{rank}(\matr{A}) = m \);
        \item There exists \( \matr{E}_i \in \mathcal{GL}(\ell) \) such that \( \matr{F}_i = \matr{F}_1\matr{E}_i \);
        \item \(\matr{F}_i^\top \matr{F}_i = \matr{I}\).
    \end{enumerate}
\end{restatable}

\begin{equation}
\label{OPP}
\min_{\matr{G}_i \in \mathcal{O}(\ell)} \quad \sum_{i=1}^{c} \| \matr{A}_i \matr{G}_i - \matr{A}_1 \matr{O} \|_\mathrm{F}^2.
\tag{OPP}
\end{equation}

\begin{definition} \label{def:orth_concordance} (Orthogonal Concordance)\\ 
The orthogonal change-of-basis matrices \(\matr{G}_i \in \mathcal{O}(\ell)\), for \( i \in [c] \), satisfy orthogonal concordance if 
\begin{equation} 
\matr{F}_1 \matr{G}_1 = \matr{F}_2 \matr{G}_2 = \dots = \matr{F}_c \matr{G}_c. 
\end{equation} 
\end{definition}

\begin{theorem}
    \label{theorem: orthogonal concordance}
    (Orthogonal Concordance of ODC)\\
    Suppose that we observe matrices \(\matr{A}_i = \matr{A}\matr{F}_i\), \( i \in [c] \), with \(\matr{A}\in \mathbb{R}^{a\times m}\) and \(\matr{F}_i\in\mathbb{R}^{m\times \ell}\). For any orthogonal matrix \(\matr{O} \in \mathcal{O}(\ell)\), let \(\matr{A}_i^\top \matr{A}_1 \matr{O} = \matr{U}_i \matr{\Sigma}_i \matr{V}_i^\top \) denote the SVD. Under \textbf{Assumption}~\ref{assumptions: ODC}, for each \( i \in [c] \), the solution \( \matr{G}_i^* = \matr{U}_i \matr{V}_i^\top\) to the Orthogonal Procrustes Problem:
    \begin{equation}
    \min_{\matr{G}_i \in \mathcal{O}(\ell)} \quad \sum_{i=1}^{c} \| \matr{A}_i \matr{G}_i - \matr{A}_1 \matr{O} \|_\mathrm{F}^2,
    \end{equation}
    satisfies
    \begin{equation}
    \matr{G}_i^* = \matr{E}_i^\top \matr{O},
    \end{equation}
    thereby guaranteeing orthogonal concordance~(\textbf{Definition}~\ref{def:orth_concordance}).
\end{theorem}

\begin{proof}
We prove for all \(i \in [c]\). Given \eqref{OPP}, we can write:
\begin{align*}
    \| \matr{A}_i \matr{G}_i - \matr{A}_1 \matr{O} \|_\mathrm{F}^2 
    &= \tr\left((\matr{A}_i \matr{G}_i - \matr{A}_1 \matr{O})^\top (\matr{A}_i \matr{G}_i - \matr{A}_1 \matr{O})\right) \\
    &= \| \matr{A}_i \|_\mathrm{F}^2 + \| \matr{A}_1 \|_\mathrm{F}^2 - 2\tr(\matr{G}_i^\top \matr{A}_i^\top \matr{A}_1 \matr{O}),
\end{align*}
where \(\tr(\cdot)\) denotes the matrix trace. Minimizing \(\| \matr{A}_i \matr{G}_i - \matr{A}_1 \matr{O} \|_\mathrm{F}^2\) for each \(\matr{G}_i\) individually is equivalent to minimizing \(\sum_{i=1}^c \| \matr{A}_i \matr{G}_i - \matr{A}_1 \matr{O} \|_\mathrm{F}^2\) for all \(\matr{G}_i\). Therefore, solving \eqref{OPP} is equivalent to solving:
\begin{equation}
\label{eq: GOPP_trace}
\begin{aligned}
\max_{\matr{G}_i \in \mathcal{O}(\ell)} \quad & \tr(\matr{G}_i^\top \matr{A}_i^\top \matr{A}_1 \matr{O}),
\end{aligned}
\end{equation}
for each \(i \in [c]\). 

From condition 3) of \textbf{Assumption}~\ref{assumptions: ODC}, we have:
\begin{equation}
\matr{A}_1 \matr{O} = \matr{A} \matr{F}_1 \matr{O} = \matr{A} \matr{F}_i \matr{E}_{i}^\top \matr{O} = \matr{A}_i \matr{E}_{i}^\top \matr{O}.
\end{equation}
Substitute this into Problem \eqref{eq: GOPP_trace}, and let \(\matr{A}_i^\top \matr{A}_i = \matr{Q}_i \matr{\Lambda}_i \matr{Q}_i^\top\) be the eigenvalue decomposition. We have:
\begin{align*}
    \tr(\matr{G}_i^\top \matr{A}_i^\top \matr{A}_1 \matr{O}) &= \tr(\matr{G}_i^\top \matr{A}_i^\top \matr{A}_i \matr{E}_{i}^\top \matr{O} ) \\
                                             &= \tr(\matr{G}_i^\top \matr{Q}_i \matr{\Lambda}_i \matr{Q}_i^\top \matr{E}_{i}^\top \matr{O} ) \\
                                             &= \tr(\matr{Q}_i^\top \matr{E}_{i}^\top \matr{O} \matr{G}_i^\top \matr{Q}_i \matr{\Lambda}_i ) \\
                                             &= \tr(\matr{W}'_i \matr{\Lambda}_i) \\
                                             &= \sum_{s=1}^{\ell} w'_{i,(s,s)} \lambda_{i,(s,s)},
                                                 \stepcounter{equation}\tag{\theequation}\label{eq: procrustes_ED}
\end{align*}
where \(\matr{W}'_i = \matr{Q}_i^\top \matr{E}_{i}^\top \matr{O} \matr{G}_i^\top \matr{Q}_i\), and \(w'_{i,(s,t)}, \lambda_{i,(s,t)}\) denote the \((s,t)\)-th elements of matrices \(\matr{W}'_i\) and \(\matr{\Lambda}_i\), respectively. Since \(\matr{W}'_i \in \mathcal{O}(\ell)\), \(w'_{i,(s,t)} \leq 1\) for all \(s, t\). Thus, the sum in \eqref{eq: procrustes_ED} is maximized when \(\matr{W}'_i = \matr{I}\), which gives:
\begin{align*}
    \matr{G}_i^* &= \matr{Q}_i \matr{Q}_i^\top \matr{E}_{i}^\top \matr{O} \\
    \matr{G}_i^* &= \matr{E}_{i}^\top \matr{O},
\end{align*}
and therefore, we have
\begin{equation}
\matr{F}_1 \matr{G}_1^* = \cdots = \matr{F}_c \matr{G}_c^*,
\end{equation}
which completes the proof.
\end{proof}

\end{document}
